% Options for packages loaded elsewhere
\PassOptionsToPackage{unicode}{hyperref}
\PassOptionsToPackage{hyphens}{url}
\documentclass[
]{article}
\usepackage{xcolor}
\usepackage{amsmath,amssymb}
\setcounter{secnumdepth}{-\maxdimen} % remove section numbering
\usepackage{iftex}
\ifPDFTeX
  \usepackage[T1]{fontenc}
  \usepackage[utf8]{inputenc}
  \usepackage{textcomp} % provide euro and other symbols
\else % if luatex or xetex
  \usepackage{unicode-math} % this also loads fontspec
  \defaultfontfeatures{Scale=MatchLowercase}
  \defaultfontfeatures[\rmfamily]{Ligatures=TeX,Scale=1}
\fi
\usepackage{lmodern}
\ifPDFTeX\else
  % xetex/luatex font selection
\fi
% Use upquote if available, for straight quotes in verbatim environments
\IfFileExists{upquote.sty}{\usepackage{upquote}}{}
\IfFileExists{microtype.sty}{% use microtype if available
  \usepackage[]{microtype}
  \UseMicrotypeSet[protrusion]{basicmath} % disable protrusion for tt fonts
}{}
\makeatletter
\@ifundefined{KOMAClassName}{% if non-KOMA class
  \IfFileExists{parskip.sty}{%
    \usepackage{parskip}
  }{% else
    \setlength{\parindent}{0pt}
    \setlength{\parskip}{6pt plus 2pt minus 1pt}}
}{% if KOMA class
  \KOMAoptions{parskip=half}}
\makeatother
\usepackage{longtable,booktabs,array}
\usepackage{caption}
\captionsetup[table]{skip=6pt}
\newcounter{none} % for unnumbered tables
\usepackage{calc} % for calculating minipage widths
% Correct order of tables after \paragraph or \subparagraph
\usepackage{etoolbox}
\makeatletter
\patchcmd\longtable{\par}{\if@noskipsec\mbox{}\fi\par}{}{}
\makeatother
% Allow footnotes in longtable head/foot
\IfFileExists{footnotehyper.sty}{\usepackage{footnotehyper}}{\usepackage{footnote}}
\makesavenoteenv{longtable}
\setlength{\emergencystretch}{3em} % prevent overfull lines
\providecommand{\tightlist}{%
  \setlength{\itemsep}{0pt}\setlength{\parskip}{0pt}}
\usepackage{bookmark}
\IfFileExists{xurl.sty}{\usepackage{xurl}}{} % add URL line breaks if available
\urlstyle{same}
% fallback for those not using the hyperref driver hyperxmp:
\makeatletter
\@ifundefined{xmpquote}{\newcommand{\xmpquote}[1]{#1}}{}
\makeatother
\hypersetup{
  hidelinks,
  pdfcreator={LaTeX via pandoc}}

\author{}
\date{}

\begin{document}

\section{Gate E2 report --- run\_v1.22\_r1 (Paper IV
v1.22-r1)}\label{gate-e2-report--run_v122_r1-paper-iv-v122-r1}

Verdict: \textbf{PASS}. Evidence: \texttt{20\_EVIDENCE/E2/}. Auditor:
Claude Opus 5.5, same session as run\_v1.2\_r1/run\_v1.21\_r1; no Lean
executed.

\begin{center}\rule{0.5\linewidth}{0.5pt}\end{center}

\section{E2 --- Focused mathematics revalidation (v1.22 against
v1.21)}\label{e2--focused-mathematics-revalidation-v122-against-v121}

Method. I compared v1.21 → v1.22 with my own unified diffs
(\texttt{auditor\_diff\_en.diff} and \texttt{auditor\_diff\_es.diff}). I
then rederived the new C.3 text by hand, block by block. For
definitions, I compared it with the pinned Mathlib source
\texttt{Mathlib/Combinatorics/SimpleGraph/Regularity/Bound.lean} (git
rev 8f9d9cff, sha256 c6395362\ldots; read-only). For the formal
statements, I compared it with \texttt{E18/Numeric.lean} and
\texttt{E19/Main.lean} in the frozen cut. The arithmetic was checked
exactly (\texttt{20\_EVIDENCE/E5/e5\_v122\_c3.py}: 14 checks, each with
a rejected negative control once the negfix and supplement files are
counted; see E5). I read the author response
(\texttt{RESPONSE\_TO\_REVALIDATION\_v1.22.md}) only after reaching
these conclusions. No Lean was run (mandate §2).

\subsection{Diff scope (v1.21 → v1.22)}\label{diff-scope-v121--v122}

\begin{itemize}
\tightlist
\item
  Status paragraphs (title page and §7 "Audit status"): no mathematics.
\item
  §5.1 heading and the A.2 row: formatting only.
\item
  A.2 scope paragraph: now refers to C.3.
\item
  C.2 now includes the R-02 clarification sentence.
\item
  C.3: new definitions of I, H, t\_j and B, plus three new paragraphs
  ("The regularity iterate at the chosen accuracy", "Absorbing the five
  terms" and "Converting the iterate to a tower").
\item
  X-12, X-15 and NEW-02 fixes: presentation only.
\item
  No theorem statement, hypothesis, final constant, equation tag or Lean
  block changed (E6\_PARITY.json: tags and Lean blocks are equal to
  v1.21; no display removed).
\end{itemize}

\subsection{C.3 --- the six mandated
checks}\label{c3--the-six-mandated-checks}

{\def\LTcaptype{none} % do not increment counter
\begin{longtable}[]{@{}llll@{}}
\toprule\noalign{}
\# & Check & Rederivation & Result \\
\midrule\noalign{}
\endhead
\bottomrule\noalign{}
\endlastfoot
1 & I(ε,k₀), H(ε), t₀, t\_\{j+1\} and B against
\texttt{SzemerediRegularity.initialBound}, \texttt{stepBound} and
\texttt{bound} & Mathlib (Bound.lean):
\texttt{stepBound\ n\ =\ n\ *\ 4\ \^{}\ n} (l.42);
\texttt{initialBound\ ε\ l\ =\ max\ 7\ (max\ l\ (⌊log\ (100/ε\^{}5)\ /\ log\ 4⌋₊\ +\ 1))}
(l.167);
\texttt{bound\ ε\ l\ =\ stepBound\^{}{[}⌊4/ε\^{}5⌋₊{]}\ (initialBound\ ε\ l)\ *\ 16\ \^{}\ (stepBound\^{}{[}⌊4/ε\^{}5⌋₊{]}\ (initialBound\ ε\ l))}
(l.190). In the manuscript, I = max\{7, k₀, ⌊log(100/ε⁵)/log 4⌋₊ + 1\},
H = ⌊4/ε⁵⌋₊, t\_\{j+1\} = t\_j 4\^{}\{t\_j\} and B = t\_H
16\^{}\{t\_H\}. & \textbf{Match} (the order inside max is immaterial) \\
2 & At η₀: I = k₀, and H equals h of §6.3
(\texttt{E18.Numeric.initialBound\_eq}, \texttt{floor\_iter\_eq},
\texttt{BE\_eta0}) & ε = δ/8, with 1/δ = 2208(9·10¹⁹)²¹ exactly. 100/ε⁵
≤ 4⁴⁰⁰⁰ (exact), so the floor term is ≤ 4001 \textless{} k₀ = 3·10¹⁹+1,
and I = k₀. 4/ε⁵ is an integer equal to 4·8⁵·2208⁵·(9·10¹⁹)¹⁰⁵ = h,
using 9·10¹⁹ = 4500·2·10¹⁶ (E5 C3a--C3c). The E18 statements have the
same content. & \textbf{Correct} \\
3 & The lower bound after two iterations dominates G₁, G₂, G₃; no claim
that two steps suffice & x·4\^{}x ≥ 2\^{}x and the iterates are
nondecreasing, so R = t\_h ≥ t₂ ≥ 2\^{}\{t₁\} ≥ 2\^{}\{2\^{}\{k₀\}\} ≥
2\^{}\{2\^{}\{59517\}\} ≥ G\_i. This uses Table C.1: G₃ ≤ 2\^{}\{U+67\}
≤ 2\^{}\{2\^{}\{59517\}\}, and h ≥ 2. The text says explicitly that this
is "a lower bound for the regularity iterate, not an upper bound on the
number of iterations needed by regularity". & \textbf{Correct} (C3g,
C3h) \\
4 & Substitution into the five terms of (C.1), including the ceilings
and the final unit (\texttt{E18.Numeric.NfarE\_eta0\_le}) & k₀ is exact.
⌈B/δ⌉ = αB exactly (B is an integer; α = 2208(9·10¹⁹)²¹).
⌈(12+10D\_*)/z⌉ = G₃ exactly, with z = 1/(2·10¹⁷). The fourth term is
⌈4B⁵G₁(a₃+a₄)/(a₃a₄)⌉, using 1/γ = G₁. Exactly, (a₃+a₄)/(a₃a₄) ≤
2(9·10¹⁹)⁶ with slack ≥ 10⁻¹⁰(9·10¹⁹)⁶, so 4G₁B⁵ times the slack is far
above 1 and the ceiling is absorbed into βG₁B⁵, β = 8(9·10¹⁹)⁶ (C3d).
⌈50(1+2C)/ξ⌉ = 10¹⁸(1+2G₂) exactly, since 50/ξ = 10¹⁸ at ξ = η₀/2 (C3e).
The leading 1 of N\_E and the max ≤ sum give the displayed N\_far(η₀) ≤
k₀ + αB + G₃ + βG₁B⁵ + 10¹⁸(1+2G₂) + 1. The Lean statement
\texttt{NfarE\_eta0\_le} is literally this sum. The rational bound for
the fourth term (the "sharper direct substitution") is correct. &
\textbf{Correct} \\
5 & B⁶/4 + 7B² ≤ B⁶ and the induction 4t\_j ≤ T(j+5) give T(h+7)
(\texttt{E19.absorb\_abstract}, \texttt{four\_mul\_Tnum\_le},
\texttt{B6\_le\_two\_pow\_two\_pow}, \texttt{NfarE\_eta0\_le\_tower}) &
k₀, α, β and 10¹⁸ are all ≤ 2²⁰⁰⁰ ≤ R (C3f). 16\^{}R ≥ 4R, so B ≥ 4R²
and B ≥ 4. Then βG₁B⁵ ≤ R²B⁵ ≤ B⁶/4. The other terms are ≤ B², B², B²,
3B² and B², i.e. 7B² in total, and 256B² ≤ B⁶, so the sum is ≤
(1/4+7/256)B⁶ ≤ B⁶ (C3i, C3j). The tower step: 4k₀ ≤ 2⁶⁷ ≤ T(5), and
4(x4\^{}x) ≤ 2\^{}\{4x\}, so 4t\_j ≤ T(j+5) (C3k, C3l). With q = T(h+5):
4R ≤ q and q ≥ 6. B⁶ = R⁶2\^{}\{24R\} ≤ 2\^{}\{30R\} ≤ 2\^{}\{8q\}
\ldots{} ≤ 2\^{}\{2\^{}q\} = T(h+7), using 30R ≤ 8q ≤ 2\^{}q (C3m, C3n).
This matches the E19 hypotheses: 2²⁰⁰⁰ ≤ T and g\_i ≤ T in
\texttt{absorb\_abstract}; 8 ≤ T and 4T ≤ t in
\texttt{B6\_le\_two\_pow\_two\_pow}. & \textbf{Correct} \\
6 & The A.2 pointer now goes to C.3; the account can be reconstructed
without delegating to theorem names & A.2 (EN p40 / ES p40) now reads
"Revision 1.22 defines the regularity recurrence in C.3 \ldots". The
dangling "recurrence in §6.3" is gone. Every step above can be followed
from the printed text and Table C.1 alone; the Lean names serve only as
cross-references. & \textbf{Correct} \\
\end{longtable}
}

\textbf{A2-4 → ACCEPTABLE\_SUMMARY. X-05 → RESOLVED.} The explicit
values T(h+7) and T(h+8) of §6.3 are now derived in the text.

\subsection{Other A.2 rows (inherited)}\label{other-a2-rows-inherited}

The diff confirms that the text of A2-1, A2-2, A2-3, A2-5, A2-6 and A2-7
is unaltered, apart from the R-02 sentence inside A2-3. They therefore
inherit their run\_v1.21\_r1 verdict of \textbf{ACCEPTABLE\_SUMMARY}.
The inheritance rests on identity: unchanged text in an identified
target and an unchanged Lean source cut (E0: 615/615).

\subsection{C.2 --- R-02 and R-01}\label{c2--r-02-and-r-01}

\begin{itemize}
\tightlist
\item
  R-02: the new sentence, "here the sum includes every active profile
  containing both classes of q, not one selected pair per profile" (EN
  p47 / ES p48), matches \texttt{PatternTransfer.servingT}. \textbf{R-02
  RESOLVED.}
\item
  R-01: the second-moment inputs B\_q u²a₃² ≤ 11δt² and 23δt² are still
  cited, not derived. As the mandate requests, this is kept as a
  declared expository limit (\textbf{OBSERVATION, KEPT}). I found no new
  reason to raise it.
\end{itemize}

\subsection{E2 verdict}\label{e2-verdict}

\textbf{PASS.} All seven A.2 rows are ACCEPTABLE\_SUMMARY. No
mathematical error was found. The scope remains a summary-level written
proof checked against frozen formal statements; it is not a replacement
for human referee review.

\end{document}
